\chapter{Anexos: Archivos Fuente y Recursos Adjuntos}
\label{sec:anexos}

Con el propósito de mantener la concisión del documento y facilitar la inspección directa del código, los archivos
fuente del proyecto se han incorporado como adjuntos embebidos en este documento PDF.

Los lectores de documentos PDF compatibles (tales como Adobe Acrobat Reader, Okular, Evince o los visores integrados
de navegadores modernos) permiten extraer, guardar y visualizar directamente los archivos haciendo clic sobre las
anotaciones de la tabla~\ref{tab:anexos_adjuntos} o accediendo al panel lateral de archivos adjuntos del visor.

\begin{table}[!ht]
  \centering
  \small
  \caption{Listado de archivos fuente y recursos técnicos embebidos en el documento.}
  \label{tab:anexos_adjuntos}
  \begin{tabular}{c l p{6.4cm} c}
    \toprule
    \textbf{Elemento} & \textbf{Archivo} & \textbf{Descripción Técnica} & \textbf{Adjunto} \\
    \midrule
    1 & \texttt{modbus\_registers.h} & Definición del mapa de registros y máscaras de bits & \attachfile[description={Archivo de cabecera modbus-registers.h}]{sources/include/modbus_registers.h} \\
    2 & \texttt{main.c} & Firmware de la central en FreeRTOS sobre STM32F103 & \attachfile[description={Codigo fuente principal main.c}]{sources/src/main.c} \\
    3 & \texttt{FreeRTOSConfig.h} & Configuración del núcleo de tiempo real FreeRTOS & \attachfile[description={Configuracion FreeRTOS}]{sources/include/FreeRTOSConfig.h} \\
    4 & \texttt{index.html} & Aplicación web SCADA y Maestro Modbus RTU & \attachfile[description={Aplicacion web index.html}]{sources/index.html} \\
    \bottomrule
  \end{tabular}
\end{table}

\section{Estructura de los Recursos Adjuntos}
\label{sec:descripcion_recursos}

A continuación se resume el alcance de cada uno de los archivos embebidos:
\begin{itemize}
  \item \texttt{modbus\_registers.h:} Cabecera en lenguaje C que tipifica los 20 registros del vector global compartido,
  las constantes de tiempo predeterminadas y las banderas de diagnóstico vial.
  \item \texttt{main.c:} Implementación completa del firmware embebido, comprendiendo la inicialización de hardware, el
  manejador de interrupciones de USART1, la máquina de estados (\texttt{vStateMachineTask}), el muestreo analógico
  (\texttt{vReadInputsTask}) y el procesador de comandos Modbus (\texttt{vSerialCommTask}).
  \item \texttt{index.html:} Aplicación autónoma desarrollada en HTML5, CSS y JavaScript que implementa la lógica de
  maestro Modbus RTU, comunicándose con el hardware mediante la interfaz Web Serial API o en modo de simulación.
  \item \texttt{FreeRTOSConfig.h:} Parámetros del planificador apropiativo, configuración de la frecuencia del tick a
  $1\unit{\kilo\hertz}$, asignación de memoria heap y prioridades de interrupción de FreeRTOS.
\end{itemize}
